\documentclass[conference]{IEEEtran}
\usepackage{cite}
\usepackage{amsmath,amssymb,amsfonts}
\usepackage{algorithmic}
\usepackage{graphicx}
\usepackage{textcomp}
\usepackage{xcolor}
\usepackage{hyperref}

\begin{document}

\title{Security Analysis of IoT Discovery and Collaboration Protocols: A DDoS Attack Scenario}

\author{\IEEEauthorblockN{Hamza Halladj}
\IEEEauthorblockA{\textit{Department of Computer Science} \\
\textit{[University Name]}\\
[City, Country] \\
email@email.com}
}

\maketitle

\begin{abstract}
This proposal outlines a methodology for evaluating the resilience of a distributed Mobile IoT discovery and collaboration protocol against Denial of Service (DoS) attacks. We propose a simulation-based approach using NS-3 to model an internal threat where a compromised mobile node initiates a high-rate UDP flood targeting a central collaboration access point. The study aims to quantify the impact on network throughput, packet loss, and channel saturation during critical collaboration phases.
\end{abstract}

\begin{IEEEkeywords}
IoT, NS-3, DDoS, UDP Flood, Network Security, Wireless Simulation
\end{IEEEkeywords}

\section{Introduction}
As Internet of Things (IoT) devices increasingly rely on ad-hoc discovery and peer-to-peer collaboration, their susceptibility to internal network attacks grows. The protocol under study involves a two-phase operation: a UDP-based Discovery Phase followed by a TCP-based Collaboration Phase. This paper proposes a scenario to test the robustness of the TCP collaboration phase against a Distributed Denial of Service (DDoS) attack launched by an insider threat.

\section{Threat Model}
We define the threat model based on the capabilities of a compromised internal node:
\begin{itemize}
    \item \textbf{Attacker Type}: Insider Threat / Compromised Mobile Node.
    \item \textbf{Access Level}: The attacker possesses valid network credentials (WPA2 keys, IP address) and is fully authenticated within the simulation environment.
    \item \textbf{Knowledge}: The attacker is aware of the protocol timeline, specifically targeting the \textit{Collaboration Phase} ($t=25s$ to $t=95s$) to maximize disruption.
    \item \textbf{Attack Vector}: UDP Flood. The attacker utilizes the connectionless nature of UDP to generate a high volume of traffic without waiting for acknowledgments.
\end{itemize}

\section{Proposed Attack Methodology}

\subsection{Simulation Setup}
The attack will be implemented within the NS-3 simulation environment (version 3.45) using the existing topological setup of fixed Access Points (APs) and mobile IoT nodes.

\subsection{Attacker Configuration}
\begin{itemize}
    \item \textbf{Source}: Node \texttt{Mobile-0} (Compromised Agent).
    \item \textbf{Target}: Node \texttt{Fixed-0} (Cluster Head/AP).
    \item \textbf{Timing}: The attack commences at $t=25.0s$, synchronized with the start of the Collaboration Phase.
\end{itemize}

\subsection{Traffic Profile}
The malicious traffic is generated using an \texttt{OnOffApplication} with a \texttt{UdpSocketFactory}. The traffic parameters are chosen to saturate the 802.11n wireless channel:

\begin{table}[htbp]
\caption{Traffic Parameters Comparison}
\begin{center}
\begin{tabular}{|l|c|c|}
\hline
\textbf{Parameter} & \textbf{Legitimate (TCP)} & \textbf{Malicious (UDP)} \\
\hline
Protocol & TCP & UDP \\
\hline
Data Rate & $\approx$ 8 kbps & 50 Mbps \\
\hline
Packet Size & 1024 bytes & 1024 bytes \\
\hline
Duty Cycle & Intermittent & 100\% (Continuous) \\
\hline
\end{tabular}
\label{tab1}
\end{center}
\end{table}

\section{Expected Impact}
We hypothesize the following impacts on the network performance:

\subsection{Channel Saturation}
The high-rate UDP stream is expected to dominate the shared wireless medium. Due to the CSMA/CA mechanism of 802.11, legitimate nodes attempting to send TCP SYN or ACK packets will persistently sense the medium as busy, leading to exponential backoff and eventual packet drops.

\subsection{Throughput Degradation}
The ``Goodput'' of legitimate collaboration links (e.g., \texttt{Mobile-1} $\to$ \texttt{Fixed-0}) is expected to degrade significantly, approaching zero as the attack intensity increases.

\subsection{Resource Exhaustion}
The receiving buffer of the target AP (\texttt{Fixed-0}) may overflow, causing tail-drop of incoming legitimate packets.

\section{Verification & Metrics}
The success of the attack will be verified using PCAP trace analysis:
\begin{enumerate}
    \item \textbf{Traffic Volume}: Visualization of I/O graphs showing UDP dominance.
    \item \textbf{Retransmissions}: Analysis of \texttt{tcp.analysis.retransmission} fields for legitimate flows.
    \item \textbf{Packet Loss Ratio}: Calculation of lost packets versus total transmitted packets for non-malicious nodes.
\end{enumerate}

\section{Conclusion}
This proposal sets the stage for a rigorous evaluation of IoT protocol security. By simulating a realistic UDP flood attack, we aim to identify critical vulnerabilities in the current collaboration implementation and propose future mitigation strategies such as rate limiting or traffic anomaly detection.

\end{document}
